\documentclass[11pt]{article}
\usepackage[margin=1in]{geometry}
\usepackage{amsmath,amssymb,booktabs}
\usepackage{amsthm}
\usepackage[breakable]{tcolorbox}

\theoremstyle{definition}
\newtheorem{proposition}{Proposition}

\tcolorboxenvironment{proposition}{
    breakable,
    colback=white,
    colframe=black!50,
    boxrule=0.4pt,
    sharp corners
}

\newtcolorbox{proofbox}{
    breakable,
    colback=white,
    colframe=black!50,
    boxrule=0.4pt,
    sharp corners
}
% Submission schedule
% Nov 23: Part 1, up to 3 pages, including Pyomo output for one family S instance.
% Dec 4:  Part 2, up to 3 pages, ending with the plan for Part 3.
% Dec 11: Final report, up to 12 pages plus appendices, and the code repository.
%
% For each experiment, write down the prediction and reasoning before
% running it. In the report, put the prediction next to the observed result.

\title{
    \large IEOR 262A: Mathematical Programming I\\
    Fall 2026 Group Project\\[0.7em]
    \LARGE Locating Capacitated Facilities: Bounds, Duality, and Branch and Bound
}
\author{
    \textbf{Team 3}\\
    Kenny Wongchamcharoen, Ashlee Liu, Helen Chang, Charmaine Yuen\\[0.4em]
    \normalsize Extension track: Track C (modular capacity)
}
\date{}

\begin{document}
\maketitle

\section{Base model}

Opening another facility can reduce shipping costs, but the savings must
justify its fixed cost. We model this trade-off first. We then examine what the LP relaxation allows a fractionally open facility to ship. The original problem is written below.

Let $I$ be the set of candidate sites and $J$ the set of customers.
Opening site $i$ costs $f_i$ and provides capacity $q_i$. Customer $j$
has demand $d_j$, and shipping one unit from $i$ to $j$ costs $c_{ij}$.
We choose which sites to open ($y_i$) and how much to ship from each
site to each customer ($x_{ij}$).

\begin{equation}
\tag{CFL}
\begin{aligned}
\min_{x,y}\quad
    & \sum_{i\in I} f_i y_i
      + \sum_{i\in I}\sum_{j\in J} c_{ij}x_{ij} \\
\text{s.t.}\quad
    & \sum_{i\in I} x_{ij} = d_j
    && j\in J, \\
    & \sum_{j\in J} x_{ij} \le q_i y_i
    && i\in I, \\
    & x_{ij}\ge 0,\quad y_i\in\{0,1\}
    && i\in I,\ j\in J.
\end{aligned}
\end{equation}

\subsection*{(a) Extension and a stronger formulation}
% Name the chosen track: A (single sourcing), B (siting rules), or
% C (modular capacity). Describe the operational decision in one sentence.
% Then define any new variables and write the changes to (CFL).
We choose Track C (modular capacity). At each candidate site $i\in I$,
we may build a small facility, build a large facility, or leave the
site closed. A small facility at site $i$ has capacity $q_i$ and
fixed cost $f_i$. A large facility at the same site has capacity
$q_i^{(2)}$ and fixed cost $f_i^{(2)}$.

For each $i\in I$, define the binary decision variables
\[
s_i =
\begin{cases}
1, & \text{if a small facility is built at site } i,\\
0, & \text{otherwise},
\end{cases}
\qquad
\ell_i =
\begin{cases}
1, & \text{if a large facility is built at site } i,\\
0, & \text{otherwise}.
\end{cases}
\]
We impose $s_i+\ell_i\le 1$ for each $i\in I$, so at most one
facility size may be chosen at a site. The choice $(s_i,\ell_i)=(0,0)$
leaves site $i$ closed. As in (CFL), $x_{ij}\ge 0$ denotes the
amount shipped from site $i\in I$ to customer $j\in J$. We formulate the modular capacity problem as
\[
\boxed{
\begin{aligned}
\min_{x,s,\ell}\quad
    & \sum_{i\in I} \bigl(f_i s_i + f_i^{(2)}\ell_i\bigr)
      + \sum_{i\in I}\sum_{j\in J} c_{ij}x_{ij} \\[0.6em]
\text{s.t.}\quad
    & \sum_{i\in I} x_{ij} = d_j
    && j\in J, \\
    & \sum_{j\in J} x_{ij} \le q_i s_i + q_i^{(2)}\ell_i
    && i\in I, \\
    & s_i + \ell_i \le 1
    && i\in I, \\
    & s_i,\ell_i \in \{0,1\}
    && i\in I, \\
    & x_{ij} \ge 0
    && i\in I,\ j\in J.
\end{aligned}
}
\tag{CFL-C}
\]
\paragraph{Family S results.}
Our assigned seeds are 301, 302, and 303.
% Family S has 20 sites and 80 customers. Team k uses seeds 100k+1,
% 100k+2, and 100k+3, with capacity ratios 2, 4, and 8.
% Before the assigned run, predict how enforcing binary openings will
% affect cost. Afterward, inspect the fractional y_i values to explain
% the difference. Include a short Pyomo output excerpt for one team run.
For family S with ratio 4 and seed 301, we predict that the LP relaxation will have a lower optimal cost than the integer model. We also expect a strict gap because the LP can use part of a site’s capacity and pay only the corresponding fraction of its opening cost. A site that serves only a few nearby customers can be worth opening fractionally in the LP. In the integer model, though, serving those customers from that site requires paying its entire opening cost. 

\begin{center}
\begin{tabular}{lrrrr}
\toprule
Instance & Ratio & Seed & LP objective & IP objective \\
\midrule
Setup check   & 4 & 1 & 42,533.67 & 49,226.88 \\
Team instance &   &   &           &           \\
\bottomrule
\end{tabular}
\end{center}

The setup run opened sites 3, 5, and 12. Requiring each site to be fully
open or closed raised the minimum cost by $6{,}693.21$ relative to the
LP relaxation.

\subsection*{(b) Strengthening inequalities}
% For each proposed inequality, explain why every feasible integer solution
% satisfies it. Then give a fractional solution: specify what each site
% opens and ships, check the original LP constraints, and show the violation.
% State which relaxation implies the inequality. Repeat for the extension.
\begin{proposition}
Inspired by the disaggregated formulation in Example~5.4.1 in the
lecture note, we decided to strengthen (CFL) by adding
\[
    x_{ij} \le d_j y_i, \qquad i\in I,\ j\in J.
\]
We call the resulting formulation $(\mathrm{CFL}^+)$.
\end{proposition}
The original
aggregated capacity constraint limits the total amount shipped from site $i$. Our new inequality also limits the share of customer $j$'s demand
that site $i$ can supply. For example, if site $i$ is 20\% open
in the LP relaxation ($y_i=0.2$), it can supply at most 20\% of
each customer's demand.

\begin{proofbox}
\emph{Proof.}
We show that the inequalities that we propose are valid. Let $(x,y)$ be feasible for (CFL) in Definition~5.1.4, and fix
$i\in I$ and $j\in J$. If $y_i=0$, then nonnegativity and the
capacity constraint give
\[
    0 \le x_{ij} \le \sum_{k\in J} x_{ik} \le q_i y_i = 0.
\]
If $y_i=1$, then nonnegativity and the demand constraint give
\[
    x_{ij} \le \sum_{h\in I} x_{hj} = d_j = d_j y_i.
\]
Thus $x_{ij}\le d_j y_i$ in both cases. Now, we give a fractional solution to the LP relaxation of (CFL). Consider one site with capacity $q_1=100$ and one customer with demand $d_1=20$. Set $y_1=1/5$ and $x_{11}=20$. The site is one fifth open and ships all 20 units to the customer. This point satisfies the original LP relaxation as demand is met, $0\le y_1\le 1$, and $x_{11}=20\le 100(1/5)=20$. However, the new inequality requires
\[
    x_{11} \le d_1 y_1 = 20(1/5)=4,
\]
which is clearly violated. Our example above shows that the LP relaxation of (CFL) does not imply $x_{ij}\le d_jy_i$ in general. However, the LP relaxation of $(\mathrm{CFL}^+)$ implies the inequality because it includes this constraint for every site-customer pair.
\hfill$\square$
\end{proofbox}
Now, we propose a strengthening of (CFL-C) following the same approach.
\begin{proposition}
For Track C, we strengthen (CFL-C) by adding
\[
    x_{ij} \le d_j(s_i+\ell_i), \qquad i\in I,\ j\in J.
\]
We call the resulting formulation $(\mathrm{CFL\text{-}C}^{+})$.
\end{proposition}

In the base model, $y_i$ indicates whether site $i$ is open.
In Track C, $s_i+\ell_i$ plays the same role: it equals one if
the site is built at either size and zero if it is closed. The
new inequality therefore limits the share of customer $j$'s demand
that site $i$ can supply, whichever size is chosen.

\begin{proofbox}
\emph{Proof.}
Let $(x,s,\ell)$ be feasible for (CFL-C), and fix $i\in I$ and
$j\in J$. Since $s_i$ and $\ell_i$ are binary and $s_i+\ell_i\le 1$,
their sum is either zero or one. If $s_i+\ell_i=0$, then both
variables are zero, so nonnegativity and the capacity constraint give
\[
    0 \le x_{ij} \le \sum_{k\in J} x_{ik}
    \le q_i s_i + q_i^{(2)}\ell_i = 0.
\]
If $s_i+\ell_i=1$, then nonnegativity and the demand constraint give
\[
    x_{ij} \le \sum_{h\in I} x_{hj}
    = d_j = d_j(s_i+\ell_i).
\]
Thus $x_{ij}\le d_j(s_i+\ell_i)$ in both cases.

Now consider the same one-site example, with small capacity
$q_1=100$, large capacity $q_1^{(2)}=150$, and customer demand
$d_1=20$. Set $s_1=1/5$, $\ell_1=0$, and $x_{11}=20$. The LP
opens one fifth of the small facility and ships all 20 units to
the customer. Both size variables lie in $[0,1]$, their sum is
$1/5\le 1$, and the shipment meets demand. The capacity constraint
also holds since $20\le 100(1/5)+150(0)=20$. However, the new
inequality requires
\[
    x_{11} \le d_1(s_1+\ell_1) = 20(1/5+0) = 4,
\]
which is violated by the shipment of 20 units. Hence the LP
relaxation of (CFL-C) does not imply the inequality in general.
The LP relaxation of $(\mathrm{CFL\text{-}C}^{+})$ implies it
because it includes this bound for every site--customer pair.
\hfill$\square$
\end{proofbox}
\subsection*{(c) Single-site convex hull}
% State your linear description of conv(X) and prove both inclusions.
% For the reverse inclusion, explain how a point with 0 < y < 1 can be
% written as a combination of a closed-site point and an open-site point.
% Identify which inequalities from (b) appear in the hull description.

\section{Lagrangian relaxation}

The demand equations link the sites: if one facility ships less to a
customer, another must make up the difference. We assign a multiplier to
each customer's demand equation and examine the shipping decision that
each site would make at those multipliers.

\subsection*{(a) Site subproblems and the greedy rule}
% Write the Lagrangian and collect the terms for one site. For an open site,
% derive the order in which customers receive capacity (Remark 5.3.3).
% Compare that value with the cost of keeping the site closed.

\subsection*{(b) Concavity and a supergradient}
% Prove concavity of w(lambda). Derive a supergradient from the shipments
% chosen in (a), checking its sign against your demand-multiplier convention.

\subsection*{(c) Subgradient method and computational bounds}
% Give the initial multipliers, step size, and stopping rule. For S and L,
% compare the best bound actually reached with the base LP, strengthened
% LP, and IP values. Explain how the chosen track changes the computation.

\subsection*{(d) Shipment bounds and initialization from LP duals}
% Remove x_ij <= d_j and repeat the comparison. Use Part 1(c) and
% Proposition 4.8.3 to explain the bounds. When initializing from the
% strengthened LP's demand duals, record the very first value of w as
% well as the best value reached.

\subsection*{(e) Aggregate capacity and the knapsack bound}
% Check whether total open capacity >= total demand is implied by the
% original LP. Retain it in the Lagrangian problem, derive the binary
% knapsack, solve it exactly, and explain the change in the bound.

% For the Dec 4 submission, end this part with a short plan for Part 3.
% Specify who will run each comparison and which quantities the logs must save.

\section{LP gaps, branch and cut, and scaling}

Two formulations can have the same optimal integer solution and still
require different amounts of search to prove optimality. We compare their
LP bounds, then use Gurobi's logs to see when an incumbent is found, how
the best bound changes, and how many nodes are explored.

% Baseline: Presolve=0, Cuts=0, Heuristics=0.
% Parts (b)--(e) use nine runs per instance and formulation: one shared
% baseline and eight variants. Change one setting or add one MIP start
% at a time so that each comparison has a clear interpretation.

\subsection*{(a) Integrality gaps and capacity ratios}
% Compute (z_IP - z_LP)/z_IP for the base and strengthened formulations.
% Use a small example to explain the effect of changing the capacity ratio.

\subsection*{(b) Primal heuristics and pruning}
% Compare baseline with Heuristics=0.05. Report root gaps, nodes, and
% runtime. Use Theorem 5.3.1 to explain which nodes an incumbent can prune.

\subsection*{(c) Branching rules and a MIP start}
% Compare default branching, VarBranch=2, and VarBranch=3 (Remark 5.3.5).
% Describe your own feasible starting solution, report its objective, and
% measure what happens to the tree when Gurobi receives it as a MIP start.

\subsection*{(d) Node selection}
% Compare baseline, MIPFocus=1, and MIPFocus=3. Report nodes, runtime, and
% time to first incumbent. Plot incumbent and best bound against nodes for
% one instance. Explain which child you would explore first in a dive.

\subsection*{(e) Root cuts}
% Compare Cuts=2 with Cuts=0 and ImpliedCuts=2. Read the bound at the end
% of the root node. Report the fraction of the gap closed, cut classes,
% nodes, and time. Compare the implied-bound result with the strengthened LP.
% Did cover cuts appear? Did the strongest root bound give the fastest solve?

\subsection*{(f) Scaling with the number of sites}
% Use 20, 40, 60, 80, and 100 sites, with five customers per site.
% Give the time limit and identify runs that reached it. Plot nodes and
% runtime on logarithmic scales, with 2^|I| as the enumeration benchmark.

\section*{Consistency checks}
% Give the repository test and result for each check: feasible solutions,
% valid lower bounds, and agreement across runs that proved optimality.
% If a check fails, report the case and explain what happened.

\section*{Sources and AI use}
% Cite course results by number where you use them. List other sources
% and tools with their purpose. For AI use, name the platform and model
% version and identify the sections, code, or figures it helped with.

\section*{Team contributions}
% Give each member a sentence identifying the code and text they wrote.

\end{document}
